%%======%%======%0%======%%=======%%
\vspace{5mm}
\section{Tecnologias utilizadas}


%%======%%======%0%======%%=======%%

A estrutura das páginas foi desenvolvida utilizando \textbf{HTML}, responsável pela organização e definição dos elementos que compõem cada página do sistema, criando assim um esqueleto para o projeto. A utilização de elementos semânticos permitiu estruturar o conteúdo de acordo com sua finalidade, contribuindo tanto para a organização do código quanto para os requisitos relacionados à acessibilidade.

A apresentação visual e a organização dos elementos foram desenvolvidas utilizando \textbf{CSS}. Os arquivos de estilo foram organizados de forma modular, separando diferentes responsabilidades do projeto. Essa organização permitiu estabelecer padrões visuais comuns entre as páginas, além de facilitar a implementação de diferentes disposições para dispositivos com tamanhos de tela distintos.

Para a implementação dos comportamentos interativos da interface foi utilizado \textbf{JavaScript}. A linguagem foi empregada para controlar funcionalidades que dependem de interação do usuário e para modificar dinamicamente determinados elementos da página. Dessa forma, HTML, CSS e JavaScript foram utilizados de maneira complementar, respectivamente, para estruturar, apresentar e proporcionar comportamento à interface.

O desenvolvimento e a organização dos arquivos foram realizados no \textbf{VS Code}, utilizado como ambiente principal de edição. Além do editor, foram utilizadas extensões que auxiliaram o processo de desenvolvimento, incluindo \textit{Prettier}, para formatação do código; \textit{Material Icon Theme}, para organização visual dos arquivos; \textit{Auto Rename Tag}, para facilitar a edição de elementos HTML; e \textit{Live Server}, utilizado para executar e visualizar localmente as páginas durante o desenvolvimento, como ja mencionado na Seção \ref{sec:processo}.

O controle de versão foi realizado por meio do \textbf{Git}, permitindo registrar a evolução do projeto e acompanhar as alterações realizadas durante o desenvolvimento. 

Para a avaliação técnica da interface foram utilizadas ferramentas de análise e validação. O \textbf{Google Lighthouse} foi utilizado para avaliar aspectos relacionados à acessibilidade, desempenho, boas práticas e otimização para mecanismos de busca. A validação da estrutura dos documentos HTML e das folhas de estilo foi realizada por meio dos validadores disponibilizados pelo \textbf{W3C}, permitindo verificar a conformidade dos códigos desenvolvidos com os padrões correspondentes.
